% Gesamttest «Quadratische Gleichungen» (GF 2.2b) — Quelle des PDFs (python3 scripts/build-lp-pdf.py leitprogramme/quadratische-gleichungen/).
% Musterlösung und Raster: bewertungspaket.tex. Alle Werte mit python3 nachgerechnet (scripts/lp/quadratische-gleichungen/zahlen.py, 11.10.2026).
% Kein \lt/\gt — pdfLaTeX kennt sie nicht, < und > tun es.
% Teile nach Kapiteln, alle ohne Hilfsmittel (RLP 2.2: «auch ohne Hilfsmittel»). Nach der Prüfung vom 12.10.2026 (M1, M2)
% neu kombiniert statt Kapitelmodelle wiederholt: Binom + Wurzel mit anderem Binom als Clip/Festhalten (G1), Teilen durch
% einen Klammerfaktor (G2a), ordnen + faktorisieren + Vieta (G2b), ordnen + Ergänzung mit Wurzelergebnis (G3), ordnen +
% Formel (G4), Typ und Verfahren begründen statt lösen + fremde Diskriminante prüfen (G5), quadratisches D(k) (G6),
% Parameter vor x² und in c mit zwei kritischen Werten (G7). Bewusst 30 statt «rund 20» Minuten (HOWTO §9).
\documentclass[11pt]{article}
\usepackage{../lp-druck}
\renewcommand{\lpname}{Quadratische Gleichungen}
\renewcommand{\lpteilgebiet}{2.2b}
\renewcommand{\lpbereich}{Grundlagenfach}
\renewcommand{\lpfuss}{Gesamttest · 25 Punkte}
\begin{document}
\lpkopf{Gesamttest}{%
  \vspace{4pt}\noindent\begin{tabularx}{\linewidth}{@{}X X@{}}
  Code (kein Name): \hrulefill & Punkte: \hrulefill\ / 25
  \end{tabularx}}

\begin{lpinfo}{Anleitung}
\textbf{Zeit:} rund 30 Minuten. \textbf{Hilfsmittel:} in allen drei Teilen keine — kein Taschenrechner (RLP 2.2: «auch ohne Hilfsmittel»).
Schreib zu jeder Aufgabe den \textbf{Rechenweg} auf — bewertet wird der Weg. Gib Lösungen als Lösungsmenge $\mathbb{L}$ an; trenne die Elemente mit Strichpunkt, z.\,B. $\{-2;\, 3\}$. Wurzeln, die nicht aufgehen, bleiben stehen (exakt).
Danach: Lösung scannen oder fotografieren (ohne Namen und Standort) und mit dem \emph{Bewertungspaket} bewerten lassen.
\end{lpinfo}

\lpteil{Teil A · Wurzelziehen und Faktorisieren}{Kapitel 1 und 2, Ordnen aus 4 · ohne Hilfsmittel · 7 Punkte}

\begin{lpaufgabe}{G1}{Binom und Wurzel}{3 P}
\begin{enumerate}[label=(\alph*),leftmargin=*,itemsep=1pt]
\item Löse $x^2 + 10x + 25 = 9$. Tipp: Links steht ein Binom.
\item Warum hat $x^2 + 10x + 25 = -1$ keine Lösung? Begründe in einem Satz.
\end{enumerate}
\lpplatz{6}
\end{lpaufgabe}

\begin{lpaufgabe}{G2}{Eine fremde Lösung und eine Zerlegung}{4 P}
\begin{enumerate}[label=(\alph*),leftmargin=*,itemsep=1pt]
\item Jemand löst $(x - 4)^2 = 3\,(x - 4)$ so: «beidseitig durch $(x - 4)$ teilen: $x - 4 = 3$, also $x = 7$». Welche Lösung geht dabei verloren, und warum? Löse die Gleichung richtig.
\item Löse $x^2 - x = 30$ durch Faktorisieren und kontrolliere die Lösungen mit dem Satz von Vieta.
\end{enumerate}
\lpplatz{7}
\end{lpaufgabe}

\lpteil{Teil B · Ergänzung, Formel, Verfahren}{Kapitel 3 und 4 · ohne Hilfsmittel · 11 Punkte}

\begin{lpaufgabe}{G3}{Quadratische Ergänzung}{3 P}
Löse $x\,(x + 6) = -2$ mit quadratischer Ergänzung. Gib die Lösungen exakt an.
\lpplatz{5}
\end{lpaufgabe}

\begin{lpaufgabe}{G4}{Erst ordnen, dann die Formel}{3 P}
Löse $(2x + 1)(x - 1) = 2$ mit der Mitternachtsformel.
\lpplatz{6}
\end{lpaufgabe}

\begin{lpaufgabe}{G5}{Verfahren begründen und einen Fehler finden}{5 P}
\begin{enumerate}[label=(\alph*),leftmargin=*,itemsep=1pt]
\item Bestimme bei (ii) zuerst den Typ. Nenne dann für jede Gleichung das schnellste Verfahren und begründe die Wahl in einem Satz. Nicht lösen.\par
  (i) $x^2 - 16x + 64 = 0$\quad (ii) $(x + 2)^2 = x^2 + 8$\quad (iii) $3x^2 + x - 5 = 0$\par\vspace{5.5\baselineskip}
\item Jemand rechnet bei $x^2 - 4x - 21 = 0$: «$D = 16 - 84 = -68$, also keine Lösung». Finde den Fehler, berechne $D$ richtig und gib die Lösungsmenge an.
\end{enumerate}
\lpplatz{5}
\end{lpaufgabe}

\lpteil{Teil C · Parameterdiskussion}{Kapitel 5 · ohne Hilfsmittel · 7 Punkte}

\begin{lpaufgabe}{G6}{Parameter im linearen und im konstanten Glied}{3 P}
Für welche $k \in \mathbb{R}$ hat $x^2 + kx + k = 0$ zwei, eine oder keine Lösung? Gib die Lösung an, wo es genau eine gibt.
\lpplatz{6}
\end{lpaufgabe}

\begin{lpaufgabe}{G7}{Parameter vor dem quadratischen Glied}{4 P}
Untersuche $k\,x^2 + 4x + k = 0$ für alle $k \in \mathbb{R}$: Für welche $k$ gibt es zwei, eine oder keine Lösung? Gib die Lösung an, wo es genau eine gibt.
\lpplatz{9}
\end{lpaufgabe}
\end{document}
